\documentclass[10pt,a4paper]{amsart}
\usepackage[margin=19mm]{geometry}
\usepackage{amsmath,amssymb,mathtools}
\usepackage[scr=rsfs,cal=euler]{mathalfa}
\usepackage{xfrac}
\usepackage[unicode,hidelinks,bookmarks=false]{hyperref}
\newcommand{\ie}{i.e.\ }
\newcommand{\cf}{cf.\ }
\newcommand{\resp}{respectively\ }
\newcommand{\Subsection}{\subsection}
\providecommand{\nobreakdash}{\nobreak\hbox{-}}
\setlength{\emergencystretch}{2em}
\multlinegap0pt
\usepackage{amssymb}
\usepackage[normalem]{ulem}
\newtheorem{lemma}{Lemma}
\newtheorem{theorem}{Theorem}
\newtheorem{claim}{Claim}
\newcommand{\MNaR}{M_{N,r}(\R)}
\newcommand{\R}{{\mathbb R}}
\newcommand{\mcQ}{{\mathcal Q}}
\title{Higher rank Gelfand-Kapranov-Zelevinsky fans}
\author{Rocco Chiriv\`i, Martina Costa Cesari, Xin Fang, Peter Littelmann}
\date{Source: arXiv:2604.20250v2 (CC BY 4.0)}
\begin{document}
\maketitle

\noindent\emph{Source and licence.} Higher rank Gelfand-Kapranov-Zelevinsky fans. Version arXiv:2604.20250v2 (CC BY 4.0), distributed by arXiv under the Creative Commons Attribution 4.0 International licence, http://creativecommons.org/licenses/by/4.0/, as recorded in the arXiv licence metadata for this exact version and shown on the abstract page https://arxiv.org/abs/2604.20250v2. Full licence notice: \texttt{licenses/arxiv-2604.20250-CC-BY-4.0.txt}.\par\smallskip

\noindent\textit{Editorial scope.} Counted unit: the article's theorem thm:gkz:complete (source0.tex lines 846--848). The article's complete original proof of it is delivered with it (source0.tex lines 850--914, one \texttt{proof} environment of 65 lines). No mathematics is written, repaired or re-derived: the statement and the proof are the article's own lines, in the article's own notation, and no retained line is substituted or re-flowed.  Retained uncounted as the source-local context the proof needs: lines 825--830.  Nothing was omitted from any retained span, so the package carries no omission record.  No retained line contains a citation, so the wrapper carries no bibliography and no bibliography entry.  No retained line contains a figure environment, an image-inclusion command or drawing code, so no image is delivered and the package declares no asset dependency.  Editorial formatting only, in addition: an AMS \texttt{amsart} preamble that restates the article's own declarations and macros used by the retained lines, namely the environments \texttt{lemma}, \texttt{theorem}, \texttt{claim} on separate counters, together with the standard packages those lines need.  The retained environments print in this document as Lemma 1, Theorem 1, Claim 1 in order of appearance, whereas the article prints the same items with the numbering of its own full document.  The complete native source of the article as distributed in the arXiv e-print for this exact version is bundled unchanged as \texttt{sources/198918/source0.tex}.
\begin{lemma}\label{lemma:max:in:polytope}
Let  $R\subseteq \mathbb R^{N}$ be a polytope. There exists a unique point 
$\mathbf p_R\in R$ such that $\mathbf p_R>\mathbf p'$ (with respect to the lexicographic order) 
for all $\mathbf p'\in R\setminus\{\mathbf p_R\}$. Moreover, 
$\mathbf p_R$ is a vertex of the polytope.
\end{lemma}

\begin{theorem}\label{thm:gkz:complete}
For each $\Psi\in \MNaR$ there exists a polytopal subdivision $\mcQ$ of $(P,A)$ such that $\Psi\in C(\mcQ,N)$.
\end{theorem}

\begin{proof}
In the following we construct a piecewise linear map $g_\Psi:P\rightarrow \mathbb R^N$ with the desired properties. The corresponding piecewise linear map $g_\Psi:K(P)\rightarrow \mathbb R^N$ is obtained via linear extension.

For $\Psi\in M_{N,r}(\mathbb R)$ let $\Omega_\Psi\subseteq \mathbb R^{n-1}\oplus \mathbb R^N$ be the polytope obtained as the convex hull:
\begin{equation}\label{equation:Omega:Psi}
\Omega_\Psi:=\textrm{convex hull of\ }\{(  v,\Psi(  v))\mid   v\in A\}\subseteq \mathbb R^{n-1}\oplus \mathbb R^N.
\end{equation}
Let $p_1: \mathbb R^{n-1}\oplus \mathbb R^N\rightarrow \mathbb R^{n-1}$ be the canonical projection onto the first component. It is an affine map such that $p_1(\Omega_\Psi)=P$. For $w\in P$, the fibre of the restricted projection 
\begin{equation}\label{equation:fibre:Omega:w}
\Omega_{ {w}}:=(p_1\vert_{\Omega_\Psi})^{-1}(w) ,
\end{equation}
can be identified with a polytope in $\mathbb R^N$. By Lemma~\ref{lemma:max:in:polytope},  the map
\begin{equation}\label{definition:gPsi}
g_{\Psi}: P\rightarrow \mathbb R^N,\quad   w\mapsto \max \Omega_{ w}
\end{equation}
is well-defined.

Let $F$ be a face of the polytope $\Omega_\Psi$ and denote by $F^o$ its relative interior. We call $F$ an \emph{upper face} if for all $(w,\mathbf u)\in F$ holds: $\mathbf u=g_\Psi(  w)$. 

\begin{claim}
A face $F$ of $\Omega_\Psi$ is an upper face if and only if there exists $w\in P$ such that $(w,g_{\Psi}(w))\in F^o$.
\end{claim}

%The upper faces can be characterized as follows:
%\begin{equation}\label{eq:characterization:of:upper:faces1}
%\textit{$F$ is an upper face}\Leftrightarrow
%\textit{$\exists\, \mathbf w\in P$ such that $(\mathbf w,g_{\Psi}(\mathbf w))\in F^o$}.
%\end{equation}
\begin{proof}
If $F$ is a point, then the conditions are clearly equivalent. Suppose $\dim F\ge 1$. Note that $\dim p_1(F)\ge 1$ because the vertices of $F$ are of the form $(w_1,\Psi(w_1))$, $\ldots$, $(w_s,\Psi(w_s))$ for some $w_1,\ldots, w_s\in A$, where $s\ge 2$.

If $F$ is an upper face, then by definition, $\mathbf u=g_{\Psi}(w)$ for all $( w,\mathbf u)\in F$. 

On the other hand, let $(  w,g_{\Psi}(  w))$ be a point in the relative interior $F^o\subset F$. We argue by contradiction. Suppose there exists a point $( y, \mathbf u)\in F$, $ y\not= w$, and an element 
$\mathbf u'\in \mathbb R^N$ such that $\mathbf u'>\mathbf u$ and $(  y, \mathbf u')\in \Omega_\Psi$. 
The ray starting in $(  y,\mathbf u)$ and passing through $(  w,g_{\Psi}( w))$ 
meets a proper face $F'\subset F$ of $F$. 
Let $( u_1,\mathbf u_1),\ldots,(  u_k,\mathbf u_k)$ be the vertices of $F'$. It follows that 
$(  w,g_{\Psi}(  w))$ is a convex linear combination of these vertices and $(  y,\mathbf u)$:
\[
(  w,g_{\Psi}(  w))= \lambda (  y,\mathbf u)+\lambda_1(  u_1,\mathbf u_1)+\ldots +\lambda_k(  u_k,\mathbf u_k),
\]
where $\lambda\not=0$. But note that $\lambda (  y,\mathbf u')+\lambda_1( u_1,\mathbf u_1)+\ldots +\lambda_k(  u_k,\mathbf u_k)=:(  w,\mathbf x)$ 
is an element in $\Omega_\Psi$ with $\mathbf x>g_{\Psi}(  w)$, which is not possible. It follows: $(  y,\mathbf u)
=(  y,g_{\Psi}(  y))$ for all $(  y,\mathbf u)\in F$, $  y\not=  w$. But now one can replace $(  w,g_{\Psi}(  w))$ 
by any other element $(  y,g_{\Psi}(  y))\in F^o$, $  y\not=  w$, and the same arguments show: for any
$(  w,\mathbf u)\in F$, $\mathbf u=g_{\Psi}(  w)$, \emph{i.e.} $F$ is an upper face.
\end{proof}

\par
\noindent
\textit{Continuation of the proof of Theorem~\ref{thm:gkz:complete}.}
Let $\mathcal F$ be the union of the upper faces. The restriction $\pi\vert_{\mathcal F}:\mathcal F\rightarrow P$ of the projection is injective by construction. For the surjectivity: take $w\in P$, the point $(w,g_{\Psi}( w))$ is contained in the boundary of the polytope $\Omega_\Psi$. Hence it is a point in the relative interior $F^o$ of a unique face $F$ of $\Omega_\Psi$. By the above claim, the face is hence an upper face, and 
$  w\in \mathrm{im\,}p_1$.

By definition, the face of an upper face is an upper face, so the images of the maximal dimensional 
upper faces $\{F_j\}_{j\in J}$ define a subdivision $\{Q_j\}_{j\in J}$ of $P$.
The restriction $g_\Psi\vert_{Q_j}$ is affine linear, it is the inverse map to the (bijective) projection 
$p_1\vert_{F_j}: F_j\rightarrow Q_j$. So $g_{\Psi}$ is piecewise linear, and, by construction, its domains of linearity are 
exactly the polytopes $Q_j$, $j\in J$. It remains to add the marking of the polytope: we set 
$A_j=\{  v\in A\cap Q_j\mid  g_\Psi( v)=\Psi({v})\}$. The set $A_j$ contains by construction all the vertices of $Q_j$. 
It follows that $\mathcal Q=(Q_j,A_j)_{j\in J}$ is a polytopal subdivision of $(P, A)$ such that $g_{\Psi}=g_{\Psi,\mathcal Q}$, 
the domains of linearity are exactly the polytopes in $\mathcal Q$, and it follows from the definition of each $A_j$ that $g_{\Psi,\mathcal Q}(v)>\Psi(v)$ 
for all $v\in A\setminus A_{\mathcal Q}$, so $\Psi\in C(\mathcal Q, N)$.
\end{proof}

\end{document}
